\documentclass[a4paper]{article}
\usepackage{amsfonts}
\usepackage{amsmath}
\usepackage{amssymb}
\usepackage{graphicx}

\newcommand{\dint}{\displaystyle\int}
\newcommand{\Bgtan}{\operatorname{Bgtan}}

\begin{document}

\noindent \large Auteur: Rob Willemsen \\
\noindent \large Studierichting: Informatica \\
\noindent \large Oefeningenreeks 4A nr. 46.7 \\

\medskip

\normalsize

$\dint \dfrac{x^5-x^4+4x+4}{x^4+4} dx$ \\

\begin{tabular}{llll|ll}
$x^5$   & $-x^4$ 	& $+4x$     & $+4$      & $x^4$     & $+4$ 	\\ \cline{5-6} 
$x^5$ 	& $+0$      & $+4x$		& 			& $x$       & $1$ 	\\ \cline{1-3}
		& $-x^4$ 	& $+0x$     & $+4$		& 			& 		\\ 
		& $-x^4$    & $+0x$     & $-4$		& 			& 		\\ \cline{2-4}
	    & 			&  	        & $8$ 	    & 			& 			 			
\end{tabular}

$= \dint (x-1) dx + \dint \dfrac{8}{x^4+4} dx$ \\

$= \dfrac{1}{2}x^2 - x + \dint \dfrac{8}{x^4+4} dx$ \\

$= \dfrac{1}{2}x^2 - x + \dint \dfrac{8}{(x^2-2x+2)(x^2+2x+2)} dx$ \\

$\Rightarrow \dint \dfrac{8}{(x^2-2x+2)(x^2+2x+2)} dx$ \\

$= \dint \dfrac{Ax+B}{x^2-2x+2} dx + \dint \dfrac{Cx+D}{x^2+2x+2} dx$ \\

$C(-1+i)+D = \left. \dfrac{8}{x^2-2x+2} \right|_{x=-1+i}$ \\

$-C + D + Ci = \dfrac{8}{(-1+i)^2-2(-1+i)+2}$ \\

$-C + D + Ci = \dfrac{8}{1-2i-1+2-2i+2}$ \\

$-C + D + Ci = \dfrac{8}{4-4i} = \dfrac{2}{1-i} = \dfrac{2(1+i)}{(1-i)(1+i)} = 1+i$ \\

$\Rightarrow C = 1$ \\

$\Rightarrow -C + D = 1 \Rightarrow D = 2$ \\

$A(1+i)+B = \left. \dfrac{8}{x^2+2x+2} \right|_{x=1+i}$ \\

$A + B + Ai = \dfrac{8}{(1+i)^2+2(1+i)+2}$ \\

$A + B + Ai = \dfrac{8}{1+2i-1+2+2i+2}$ \\

$A + B + Ai = \dfrac{8}{4+4i} = \dfrac{2}{1+i} = \dfrac{2(1-i)}{(1+i)(1-i)} = 1-i$ \\

$\Rightarrow A = -1$ \\

$\Rightarrow A + B = 1 \Rightarrow B = 2$ \\

We lossen eerst de twee integralen apart op. \\

$I_1 = \dint \dfrac{-x+2}{x^2-2x+2} dx$ \\

$= -\dfrac{1}{2} \dint \dfrac{2x-4}{x^2-2x+2} dx$ \\

$= -\dfrac{1}{2} \dint \dfrac{2x-2-2}{x^2-2x+2} dx$ \\

$= -\dfrac{1}{2} \dint \dfrac{2x-2}{x^2-2x+2} dx + \dint \dfrac{1}{x^2-2x+2} dx$ \\

$= -\dfrac{1}{2} \ln|x^2-2x+2| + \dint \dfrac{1}{(x-1)^2+1} dx$ \\

$= -\dfrac{1}{2} \ln(x^2-2x+2) + \Bgtan(x-1)$ \\

$I_2 = \dint \dfrac{x+2}{x^2+2x+2} dx$ \\

$= \dfrac{1}{2} \dint \dfrac{2x+4}{x^2+2x+2} dx$ \\

$= \dfrac{1}{2} \dint \dfrac{2x+2+2}{x^2+2x+2} dx$ \\

$= \dfrac{1}{2} \dint \dfrac{2x+2}{x^2+2x+2} dx + \dint \dfrac{1}{x^2+2x+2} dx$ \\

$= \dfrac{1}{2} \ln|x^2+2x+2| + \dint \dfrac{1}{(x+1)^2+1} dx$ \\

$= \dfrac{1}{2} \ln(x^2+2x+2) + \Bgtan(x+1)$ \\

Totale integraal: \\

$= \dfrac{1}{2}x^2 - x + I_1 + I_2$ \\

$= \dfrac{1}{2}x^2 - x - \dfrac{1}{2} \ln(x^2-2x+2) + \Bgtan(x-1) + \dfrac{1}{2} \ln(x^2+2x+2) + \Bgtan(x+1) + C$ \\




\begin{thebibliography}{999}
%\bibitem{a} Referentie
\end{thebibliography}
\end{document}